\section{켤레원소와 체매장}
\label{sec:conjugates-embeddings}

\begin{learninggoal}
대수적 원소의 켤레원소를 최소다항식, 단순확장의 동형, 대수적 폐포의 자동동형이라는 세 관점으로 이해한다. 체매장이 생성원의 켤레를 어떻게 선택하는지 계산한다.
\end{learninggoal}

체 $K$의 대수적 폐포 $\overline K$를 하나 고정한다. $K$ 위에서 대수적인 원소 $\alpha\in\overline K$의 최소다항식을 $m_{\alpha,K}$라 쓰자.

\begin{definition}[켤레원소]
$K$ 위에서 대수적인 두 원소 $\alpha,\beta\in\overline K$가 같은 최소다항식을 가지면 $\alpha$와 $\beta$를 $K$ 위의 \emph{켤레원소}(conjugate elements)라 한다. $\alpha$의 모든 $K$-켤레의 집합을
\[
  \operatorname{Conj}_K(\alpha)
  :=\{\beta\in\overline K:m_{\beta,K}=m_{\alpha,K}\}
\]
로 쓴다.
\end{definition}

최소다항식은 기약이므로 $\operatorname{Conj}_K(\alpha)$는 정확히 $m_{\alpha,K}$의 서로 다른 근들의 집합이다. 중근은 한 번만 센다.

\begin{theorem}[켤레원소의 동치조건]
$\alpha,\beta\in\overline K$가 $K$ 위에서 대수적이라 하자. 다음은 동치이다.
\begin{enumerate}[label=(\roman*)]
  \item $\alpha$와 $\beta$는 $K$ 위의 켤레원소이다.
  \item 유일한 $K$-동형
  \[
    K(\alpha)\longrightarrow K(\beta),
    \qquad \alpha\longmapsto\beta
  \]
  가 존재한다.
  \item 어떤 $K$-자동동형 $\sigma\in\Aut_K(\overline K)$가 존재하여
  \[
    \sigma(\alpha)=\beta
  \]
  이다.
\end{enumerate}
\end{theorem}

\begin{proof}
(i)$\Rightarrow$(ii): 공통 최소다항식을 $m$이라 하자. 두 체는 모두 $K[X]/(m)$과 $K$-동형이므로
\[
  g(\alpha)\longmapsto g(\beta)
\]
로 주어지는 $K$-동형이 존재한다. 생성원의 상이 정해졌으므로 유일하다.

(ii)$\Rightarrow$(iii): 대수적 확장에 대한 체매장 연장 정리를 적용하면 $K(\alpha)\to K(\beta)\subseteq\overline K$를 $K$-자동동형 $\overline K\to\overline K$로 연장할 수 있다.

(iii)$\Rightarrow$(i): $m_{\alpha,K}(\alpha)=0$이고 $\sigma$가 $K$를 고정하므로
\[
  m_{\alpha,K}(\beta)
  =m_{\alpha,K}(\sigma(\alpha))
  =\sigma(m_{\alpha,K}(\alpha))=0.
\]
따라서 $\beta$는 $m_{\alpha,K}$의 근이다. 두 원소의 최소다항식은 모두 일계수 기약이고 하나가 다른 하나를 나누므로 서로 같다.
\end{proof}

\begin{corollary}
$\Aut_K(\overline K)$가 $\overline K$에 작용할 때 $\alpha$의 궤도는 켤레집합과 같다.
\[
  \Aut_K(\overline K)\cdot\alpha
  =\operatorname{Conj}_K(\alpha).
\]
또한
\[
  \#\operatorname{Conj}_K(\alpha)
  =[K(\alpha):K]_{\mathrm s}.
\]
\end{corollary}

\begin{proof}
첫 등식은 앞 정리의 (i)와 (iii)의 동치이다. 둘째 등식은 $K$-매장 $K(\alpha)\to\overline K$가 $\alpha$의 서로 다른 켤레를 선택하는 것과 일대일로 대응한다는 분리차수의 단순확장 계산이다.
\end{proof}

\begin{example}
$\alpha=\sqrt2$의 $\Q$-켤레는
\[
  \sqrt2,\quad -\sqrt2
\]
이다. 두 $\Q$-매장 $\Q(\sqrt2)\to\C$는 $\sqrt2$를 각각 이 두 원소로 보낸다.
\end{example}

\begin{example}
$\alpha=\sqrt[3]{2}$의 $\Q$-켤레는
\[
  \alpha,\quad \omega\alpha,\quad \omega^2\alpha,
  \qquad \omega=e^{2\pi i/3}
\]
이다. 따라서 $\Q(\alpha)$에서 $\C$로 가는 $\Q$-매장은 세 개이다. 그러나 $\omega\alpha$와 $\omega^2\alpha$는 실수체 $\Q(\alpha)$ 안에 들어 있지 않는다. 켤레원소가 반드시 원래 단순확장 안에 있을 필요는 없다.
\end{example}

\begin{example}[순수비분리 원소]
$K=\F_p(t^p)$, $\alpha=t$라 하자. 최소다항식은
\[
  X^p-t^p=(X-t)^p
\]
이고 서로 다른 근은 $t$ 하나뿐이다. 따라서
\[
  \operatorname{Conj}_K(t)=\{t\},
  \qquad
  [K(t):K]_{\mathrm s}=1
\]
이지만 보통 차수는 $p$이다.
\end{example}

\subsection{여러 생성원과 체매장}

$L=K(\alpha_1,\dots,\alpha_r)$가 유한 대수적 확장이라 하자. $K$-매장
\[
  \sigma:L\longrightarrow\overline K
\]
는 생성원들의 상
\[
  \sigma(\alpha_1),\dots,\sigma(\alpha_r)
\]
에 의해 유일하게 정해진다. 각 $\sigma(\alpha_i)$는 $\alpha_i$의 $K$-켤레이지만, 켤레들을 서로 독립적으로 아무렇게나 고를 수 있는 것은 아니다. 생성원들이 만족하는 모든 대수적 관계도 보존되어야 한다.

\begin{example}
$L=\Q(\sqrt2,\sqrt3)$에서 $\sqrt2$와 $\sqrt3$의 부호는 독립적으로 선택할 수 있다. 따라서 네 개의 $\Q$-매장이 존재한다. 반면 $L=\Q(\sqrt2,\sqrt6)$에서는
\[
  \sqrt6=\sqrt2\sqrt3
\]
와 같은 관계 때문에 생성원의 상들이 곱셈관계를 만족해야 한다.
\end{example}

\begin{keyidea}
켤레원소는 하나의 대수적 수가 기준체에서 보기에 가질 수 있는 모든 모습이다. 체매장은 이 모습들 가운데 하나를 선택하며, 정규확장은 그 모든 모습을 체 내부에 함께 담는 확장이다.
\end{keyidea}

\begin{checkpoint}
\begin{enumerate}[label=(\alph*)]
  \item $\sqrt2+\sqrt3$의 $\Q$-켤레를 모두 구하라.
  \item $i$와 $-i$가 $\R$ 위에서는 켤레인지 판정하라.
  \item 순수비분리 원소의 켤레집합이 한 원소만 갖는 이유를 설명하라.
\end{enumerate}
\end{checkpoint}
